\documentclass[11pt,a4paper]{article}
\usepackage[T1]{fontenc}
\usepackage{lmodern,amsmath,amssymb,amsthm,microtype}
\usepackage[margin=26mm]{geometry}
\usepackage[colorlinks=true,linkcolor=blue,urlcolor=blue]{hyperref}
\newtheorem{theorem}{Theorem}[section]
\newtheorem{lemma}[theorem]{Lemma}
\newcommand{\Q}{\mathbb Q}\newcommand{\Z}{\mathbb Z}
\title{A pole-cancelling height coordinate for an elliptic quotient}
\author{}\date{17 November 2026}
\begin{document}\maketitle

\begin{abstract}
A rational coordinate cancels the pole at the elliptic identity when
the external input equals an integer quotient of point indices. Its
height controls both that quotient and the input. The exact leading
formal coefficient gives simultaneous preparation at any prescribed
finite set of good places. All estimates are uniform in the input,
the two indices and the positive coefficient shift.
\end{abstract}

\section{Height and local cancellation}

Fix $d\in\Z_{>0}$, $E_d:y^2=x^3-d^2x$, and a nontorsion point
$P\in E_d(\Q)$. Let $H=[h]P$, $U=[j]P$, where $h,j$ are
nonzero integers. Suppose that $h\mid j$, and write $n=j/h$, so
$U=[n]H$. For a nonzero rational input $X$ and a positive integer
$B$, put
\[
 \beta=X^2+B,\qquad
 Y=X^2x(U)-x(H),\qquad
 E_h=h(1:\beta:Y).
\]
The letter $Y$ denotes a rational height coordinate, not a pole
ordinate. Let $C_E$ be a constant such that
$|h(x(R))-2\widehat h(R)|\le C_E$ for every nonzero rational
point $R$ of $E_d$; such a constant exists by
\cite[Chapter VIII, Section 9]{AEC}.

\begin{lemma}\label{lem:height}
Uniformly in $X,B,h,j$ subject to these hypotheses,
\[
 E_h\ge2h(X),\qquad
 E_h\ge(n^2-1)\widehat h(H)-C_E-\tfrac12\log2.
\]
In particular
$E_h\ge(n^2-1)\widehat h(P)-C_E-\tfrac12\log2$.
\end{lemma}
\begin{proof}
For $X=a/b$ in lowest terms with $b>0$, the fraction
$\beta=(a^2+Bb^2)/b^2$ is reduced, and its positive numerator
dominates both $a^2$ and $b^2$. Hence $h(\beta)\ge2h(X)$.
The height of a projective vector dominates the height of every
ratio of its coordinates, so $E_h\ge\max(h(\beta),h(Y))$.

Since $X\ne0$, the standard sum and product inequalities give
\[
 h(x(U))\le h(Y)+h(x(H))+2h(X)+\log2.
\]
Consequently
\[
 h(Y)+2h(X)\ge2(n^2-1)\widehat h(H)-2C_E-\log2.
\]
Use $2E_h\ge h(Y)+2h(X)$ and
$\widehat h(H)=h^2\widehat h(P)\ge\widehat h(P)>0$.
\end{proof}

\begin{lemma}\label{lem:local}
For the formal parameter $t=-x/y$ at the identity of $E_d$,
\[
 n^2x([n]H)-x(H)
       =\frac{d^2(n^4-1)}5\,t^2+O(t^6)\quad\text{in }\Q((t)).
\]
Thus on the intended input $X=n$, the function $Y$ is regular
at $H=O$ and vanishes there. For $p\nmid 10dn$ and formal
$H\in E_d(\Q_p)$ one has $v_p(Y)\ge2v_p(t)$.
For $n=1$ or $-1$, the function $Y$ is identically zero.
\end{lemma}
\begin{proof}
The curve equation and the choice of $t$ give
\[
 x=t^{-2}+d^2t^2+O(t^6),\qquad
 \frac{dx}{2y}=(1-2d^2t^4+O(t^8))\,dt.
\]
The formal logarithm is therefore
$\ell(t)=t-(2d^2/5)t^5+O(t^9)$, and its inverse gives
\[
 [n]t=nt+\frac{2d^2}{5}(n^5-n)t^5+O(t^9).
\]
Substitution in $x([n]t)$ yields the displayed coefficient.
At $p\nmid10dn$, the curve and its formal group have integral
coefficients and $n$ is a unit. After cancellation of the leading
Laurent term the remaining series belongs to $t^2\Z_p[[t]]$.
\end{proof}

For even $e\ge4$, define
\[
 M=\beta^e+Y^e.
\]
This quantity is positive. On $X=n$, at every prime in
Lemma~\ref{lem:local} where $\beta$ is integral,
\[
 M\equiv\beta^e\pmod{p^{2e v_p(t)}}.
\]
Hence a finite collection of simple contact congruences involving
$M$ can be prepared by first using the limiting value $\beta^e$
and then taking $H$ sufficiently deep in the corresponding formal
groups. This local statement makes no assertion about primes
outside the prescribed collection.

\begin{thebibliography}{9}
\bibitem{AEC}J.~H.~Silverman, \emph{The Arithmetic of Elliptic Curves},
second edition, Graduate Texts in Mathematics 106, Springer, 2009,
Chapter VIII, Section 9.
\end{thebibliography}
\end{document}
